{\heiti 摘要：}
随着海洋大数据的计算引擎从岸基超算中心延伸至移动观测平台边缘，自主水下航行器在有限的机载算力与观测能量约束下，必须在继续计算、补充观测与立即行动之间做出在线权衡。在台湾以东至东海的黑潮强流区，当航行器仅携带数天前生成的粗分辨率海流图绕越障碍时，海流速度与航行器对水速度相当，选错绕行一侧将导致数小时的航时损失。现有确定性规划方法在最佳估计地图上增加计算量以加密搜索网格，但只能提升求解精度而无法补充缺失的海洋信息；现有集合规划与自适应采样方法或因样本有限而漏掉翻转决策的极端流场，或将观测耗费在与当前航路选择无关的全场重构上。为此，本文提出一种基于海洋大数据标定的合理反例检验方法，通过分离数值歧义与环境歧义来决定何时计算、何时观测、何时采取保守航路。该方法首先通过粗精网格容限消除由求解器粗糙引起的数值歧义，随后利用多年历史再分析海流标定旧粗图的无散误差模态与合理集合半径，通过基于包络定理的投影梯度上升在合理集合内搜索使所选航路变差、另一侧变好的合理反例及最坏后悔，并据此门控触发上浮观测或最坏后悔最小的保守选择。基于六年黑潮区域再分析海流数据的仿真实验表明，将航线求解计算量增加64倍仅将选错一侧的比例从21.7\%微降至19.7\%，而在包住真实海洋94\%任务两侧时间差的合理集合内，99.7\%的任务存在合理反例；按反例最坏后悔排名前20\%的任务集中了67\%的航时损失，门控后取保守策略将平均航时损失从照精算地图走的0.98小时显著降至0.70小时（降低28\%），并在上浮通信代价低于约1小时时有效触发上浮取新图。该研究揭示了机载计算与现场观测在消除不同来源不确定性上的本质边界，为受限边缘平台提供了面向决策的轻量化元推理准则。

{\heiti 关键词：}海洋大数据；自主水下航行器；航路规划；合理反例检验；元推理；不确定性量化

\vspace{1em}
\begin{center}
{\large\bfseries Plausible Counterexample Test: When an Autonomous Underwater Vehicle Should Compute and When It Should Observe Under a Coarse Current Map}
\end{center}
\vspace{0.5em}

{\bfseries Abstract:}
As the computing engine of ocean big data extends from shore-based supercomputing centers to the edge of mobile observation platforms, an autonomous underwater vehicle operating under limited onboard computing power and observation energy must make online trade-offs among continued computation, additional observation, and immediate movement. In strong current regions such as the Kuroshio east of Taiwan and the East China Sea, when a vehicle navigates around an obstacle carrying only a multi-day-old coarse-resolution current map, current speeds are comparable to the vehicle's through-water speed, and choosing the wrong side of the obstacle costs hours of transit time. Existing deterministic planning methods increase computation on the best-estimate map to refine the search grid, which only sharpens numerical precision without supplying missing oceanographic information; meanwhile, ensemble planning and adaptive sampling methods either miss decision-flipping extreme flow fields due to finite ensemble sizes or expend observation budgets reconstructing the entire flow field rather than resolving the binary route choice. To address this, we propose a plausible counterexample test calibrated by ocean big data that separates numerical ambiguity from environmental ambiguity to decide when to compute, when to observe, and when to hedge conservatively. The method first resolves numerical ambiguity caused by coarse solvers via a coarse-to-fine numerical tolerance check, then uses multi-year historical reanalysis currents to calibrate divergence-free error modes and the radius of a plausible set for stale coarse maps, and applies envelope-theorem-based projected gradient ascent within the plausible ball to search for the worst plausible counterexample that flips the preferred side and quantifies its worst regret, gating between surfacing for a new map and choosing the side with the smaller worst regret. Experiments over six years of reanalysis currents in the Kuroshio region show that increasing route-solving computation by 64 times reduces the wrong-side rate only from 21.7\% to 19.7\%, whereas within the calibrated plausible set covering 94\% of real time differences on test years, 99.7\% of missions admit a plausible counterexample; the top 20\% of missions ranked by worst plausible regret capture 67\% of total map loss, and the gated conservative response reduces mean transit loss by 28\% from 0.98 hours to 0.70 hours while appropriately triggering surfacing when communication costs less than about 1 hour. This study establishes the fundamental boundary between onboard computation and in-situ observation in removing distinct sources of uncertainty, providing a decision-oriented lightweight metareasoning criterion for resource-constrained ocean edge platforms.

{\bfseries Keywords:} ocean big data; autonomous underwater vehicle; path planning; plausible counterexample test; metareasoning; uncertainty quantification
